% Fragmento da Introdução para o template SBC (Overleaf).
% Issue #12 — colar no corpo do artigo. Atualize o link do projeto Overleaf no README.

\section{Introdução}

As métricas DORA (\textit{DevOps Research and Assessment}) tornaram-se referência
para avaliar o desempenho de entrega de software. Popularizadas por
\textit{Accelerate}~\cite{forsgren2018accelerate} e atualizadas pelos relatórios
\textit{Accelerate State of DevOps}, as quatro métricas clássicas --- frequência
de deploy, \textit{lead time for changes}, taxa de falha de mudança
(\textit{change failure rate}, CFR) e tempo de recuperação --- capturam
\textit{throughput} e estabilidade.

Medir essas métricas em repositórios \textit{open source} exige
\textbf{proxies}: o GitHub não registra ``deploy em produção'' nem ``falha em
produção''. Releases publicadas e execuções de CI/CD (GitHub Actions) aproximam
entrega e falha, mas podem distorcer conclusões. Por isso, além de calcular as
métricas, este trabalho valida heurísticas automáticas contra julgamento humano
e avalia a sensibilidade das classificações DORA a mudanças nas definições
operacionais.

O objetivo é minerar as métricas DORA a partir de repositórios populares que
usam GitHub Actions, com definições operacionais alinhadas à turma, produzir um
\textit{pipeline} reprodutível e discutir hipóteses informais \textbf{antes} da
análise dos dados.

\subsection*{Hipóteses informais}

\begin{itemize}
  \item \textbf{RQ\,01 (A):} poucos repositórios alcançarão frequência
    \textit{Elite}; a mediana tenderá a \textit{Medium}/\textit{High}.
  \item \textbf{RQ\,02 (B):} lead time (a) $>$ lead time (b), por commits
    antigos que inflam a variante por release.
  \item \textbf{RQ\,03a (C):} CFR de CI $>$ CFR de entrega (falhas de
    \textit{pipeline} $\neq$ falhas em produção).
  \item \textbf{RQ\,03b (B):} CFR (b) concentrará repositórios em
    \textit{Elite}/\textit{High}, sujeita à heurística de release corretiva.
  \item \textbf{RQ\,04 (C):} recuperação tipicamente em horas a poucos dias, com
    censura não desprezível.
  \item \textbf{RQ\,05 (B):} correlação Spearman entre frequência de deploy e CFR
    fraca ou nula (sem \textit{trade-off} forte).
  \item \textbf{RQ\,06 (A):} ao menos um fator
    (linguagem/estrelas/contribuidores/idade/tipo) diferenciará alguma métrica,
    com efeito pequeno a moderado.
  \item \textbf{RQ\,07 (C):} mudança relevante de categoria DORA entre
    definições; kappa ponderado no máximo moderado.
\end{itemize}

% Nota: adicione a entrada forsgren2018accelerate na bibliografia do template,
% ou ajuste a citação conforme o .bib do grupo.
